\section{Merge moments on paired copy spaces}
\label{sec:pair_merge_moment}

Let $Q$ and $C$ be finite orthonormal bases and $m\ge0$. The separate
copy strings give the basis $Q^m\times C^m$ of
$\mathcal H_m=\ell^2(Q)^{\otimes m}\otimes\ell^2(C)^{\otimes m}$.
Write $q=|Q|\,|C|$. The merge-moment estimate in
\cite[Lemma~6.1(3)]{OpenAI2026AreaLaw} applies directly in these coordinates.
Its homogeneous form retains the actual mass of a positive component,
as needed in the merge-moment argument of Section~7 of that source.
The construction of the physical components and the exponential deficit
interpretation require additional arguments.

\begin{definition}[The separate and simultaneous copy actions]
\label{def:pair_copy_actions}
\lean{TensorPower.pairCopyLeft, TensorPower.pairCopyRight,
  TensorPower.pairCopyBoth, TensorPower.commute_pairCopy,
  TensorPower.pairCopyBoth_eq_mul}
\leanok
For $\sigma\in S_m$ and $(x,c)\in Q^m\times C^m$, put
\begin{align}
 u_Q(\sigma)(x,c)&=(x\circ\sigma^{-1},c), &
 u_C(\sigma)(x,c)&=(x,c\circ\sigma^{-1}),\notag\\
 u_{QC}(\sigma)(x,c)&=(x\circ\sigma^{-1},c\circ\sigma^{-1}).
 \label{eq:pair_copy_actions}
\end{align}
Let $U_Q,U_C,U_{QC}$ be their permutation operators. The separate actions
commute, and $u_{QC}(\sigma)=u_Q(\sigma)u_C(\sigma)$.
Let $\Pi^Q_\lambda,\Pi^C_\mu,\Pi^{QC}_\nu$ be the corresponding central
Schur projectors, with irreducible dimensions $d_\lambda,d_\mu,d_\nu$.
\end{definition}
\begin{proof}
\leanok
The two separate actions change different coordinates of $(x,c)$.
Applying both with the same permutation gives the simultaneous action.
\end{proof}

\begin{theorem}[The normalized paired merge moment]
\label{thm:pair_merge_moment}
\lean{TensorPower.pair_merge_moment_le}
\leanok
\uses{def:pair_copy_actions,thm:schur_label_merge_moment}
Suppose $\rho\ge0$, $\Tr\rho=1$, and $\rho$ commutes with
$U_Q(\sigma)$ and $U_C(\sigma)$ for every $\sigma\in S_m$.
For every real $b\le1$,
\begin{align}
 \sum_{\lambda,\mu,\nu}
  \Re\Tr\!\left(\rho\Pi^Q_\lambda\Pi^C_\mu\Pi^{QC}_\nu\right)
  \left(\frac{d_\lambda d_\mu}{d_\nu}\right)^b
  &\le(m+1)^{q^2}.
 \label{eq:pair_merge_moment}
\end{align}
\end{theorem}
\begin{proof}
\leanok
\uses{thm:schur_label_merge_moment}
Identify a configuration with its two copy strings:
\[
 ((x_j,c_j))_{j=0}^{m-1}
   \longmapsto ((x_j)_{j=0}^{m-1},(c_j)_{j=0}^{m-1}).
\]
This bijection transports the density and all three copy actions. It
therefore transports their central projectors and preserves every trace
in \eqref{eq:pair_merge_moment}. The separate invariances are unchanged,
and the one-copy dimension is exactly $q$. The proof of
Theorem~\ref{thm:schur_label_merge_moment} uses
$d_\lambda d_\mu/d_\nu\ge1$ on every nonzero joint block. Consequently,
its bound applies for every $b\le1$, including negative $b$.
\end{proof}

\begin{theorem}[The paired merge moment with its actual trace mass]
\label{thm:pair_merge_moment_mass}
\lean{TensorPower.pair_merge_moment_le_mul_trace}
\leanok
\uses{def:pair_copy_actions,thm:pair_merge_moment}
Suppose $\rho\ge0$ and $\rho$ commutes with both separate copy actions.
No normalization or positive-mass condition is imposed. For every real
$b\le1$,
\begin{align}
 \sum_{\lambda,\mu,\nu}
  \Re\Tr\!\left(\rho\Pi^Q_\lambda\Pi^C_\mu\Pi^{QC}_\nu\right)
  \left(\frac{d_\lambda d_\mu}{d_\nu}\right)^b
  &\le(m+1)^{q^2}\Re\Tr\rho.
 \label{eq:pair_merge_moment_mass}
\end{align}
The statement includes $m=0$, empty bases and zero-trace components.
\end{theorem}
\begin{proof}
\leanok
\uses{thm:pair_merge_moment}
Set $r=\Re\Tr\rho\ge0$. If the basis is empty then $\rho=0$.
Otherwise choose the total normalization
\[
 \eta=\begin{cases}
   \rho/r,&r>0,\\
   I/\dim\mathcal H_m,&r=0.
 \end{cases}
\]
It is positive, has trace one and retains both separate invariances.
Positivity implies $\rho=0$ when $r=0$, so $r\eta=\rho$ in both cases.
Apply \eqref{eq:pair_merge_moment} to $\eta$ and multiply by $r$.
Trace linearity gives \eqref{eq:pair_merge_moment_mass}.
\end{proof}
